\documentclass[10pt,a4paper]{amsart}
\usepackage[margin=19mm]{geometry}
\usepackage{amsmath,amssymb,mathtools}
\usepackage[scr=rsfs,cal=euler]{mathalfa}
\usepackage{xfrac}
\usepackage[unicode,hidelinks,bookmarks=false]{hyperref}
\newcommand{\ie}{i.e.\ }
\newcommand{\cf}{cf.\ }
\newcommand{\resp}{respectively\ }
\newcommand{\Subsection}{\subsection}
\providecommand{\nobreakdash}{\nobreak\hbox{-}}
\setlength{\emergencystretch}{2em}
\multlinegap0pt
\setcounter{tocdepth}{1}
\usepackage{epsfig}
\usepackage{color}
\definecolor{shadecolor}{gray}{0.875}
\usepackage{amscd}
\usepackage{adjustbox}
\usepackage{multirow}
\usepackage{tikz-cd}
\usepackage{xy}
\numberwithin{equation}{section}
\usetikzlibrary{positioning}
  \def\gn#1#2{{$\href{http://groupnames.org/\#?#1}{#2}$}}
\def\gn#1#2{$#2$}  % comment this line out to get html links
\tikzset{sgplattice/.style={inner sep=1pt,norm/.style={red!50!blue},char/.style={blue!50!black},
  lin/.style={black!50}},cnj/.style={black!50,yshift=-2.5pt,left=-1pt of #1,scale=0.5,fill=white}}
\newtheorem{prop}{Proposition}[section]
\newtheorem{theorem}[prop]{Theorem}
\newtheorem{theo}[prop]{Theorem}
\newtheorem{coro}[prop]{Corollary}
\newtheorem{cond}[prop]{Condition}
\newtheorem{lemm}[prop]{Lemma}
\newtheorem{assu}[prop]{Assumption}
\newtheorem{sublemm}[prop]{Sublemma}
\newtheorem{defi}[prop]{Definition}
\newtheorem{exam}[prop]{Example}
\newtheorem{prob}[prop]{Problem}
\newtheorem{ques}[prop]{Question}
\newtheorem{conjecture}[prop]{Conjecture}
\newtheorem{hypo}[prop]{Hypothesis}
\newtheorem{rema}[prop]{Remark}
\newtheorem{rems}[prop]{Remarks}
\newtheorem{example}[prop]{Example}
\newtheorem{exer}[prop]{Exercise}
\newtheorem{nota}[prop]{Notation}
\newtheorem{probn}{Problem}
\newcommand{\marg}[1]{\footnote{#1}\marginpar[\hfill\tiny\thefootnote$\rightarrow$]{$\leftarrow$\tiny\thefootnote}}
\newcommand\commentr[1]{{\color{red} \sf [#1]}}
\newcommand\commentb[1]{{\color{blue} \sf [#1]}}
\newcommand\commentm[1]{{\color{magenta} \sf [#1]}}
\newcommand{\sho}[1]{{\color{blue} \sf $\clubsuit\clubsuit\clubsuit$ Sho: [#1]}}
\newcommand{\yuri}[1]{{\color{red} \sf $\clubsuit\clubsuit\clubsuit$ yuri: [#1]}}
\newcommand{\upperhalfplane}{{\mathbb H}}
\newcommand{\eqto}{\stackrel{\lower1.5pt\hbox{$\scriptstyle\sim\,$}}\to}
\newcommand{\eqdashto}{\stackrel{\lower1.5pt\hbox{$\scriptstyle\sim\,$}}\dashrightarrow}
\newcommand{\actsfromleft}{\mathrel{\reflectbox{$\righttoleftarrow$}}}
\newcommand{\actsfromright}{\righttoleftarrow}
\def\acts{\actsfromleft}
\def\lra{\longrightarrow}
\def\ra{\rightarrow}
\def\da{\downarrow}
\def\cA{{\mathcal A}}
\def\cB{{\mathcal B}}
\def\cBC{{\mathcal{BC}}}
\def\BC{{\mathcal{BC}}}
\def\SC{{\mathcal{SC}}}
\def\cC{{\mathcal C}}
\def\cD{{\mathcal D}}
\def\cE{{\mathcal E}}
\def\cF{{\mathcal F}}
\def\cG{{\mathcal G}}
\def\cH{{\mathcal H}}
\def\cI{{\mathcal I}}
\def\cJ{{\mathcal J}}
\def\cK{{\mathcal K}}
\def\cL{{\mathcal L}}
\def\cM{{\mathcal M}}
\def\cN{{\mathcal N}}
\def\cO{{\mathcal O}}
\def\cP{{\mathcal P}}
\def\cQ{{\mathcal Q}}
\def\cR{{\mathcal R}}
\def\cS{{\mathcal S}}
\def\cT{{\mathcal T}}
\def\cU{{\mathcal U}}
\def\cV{{\mathcal V}}
\def\cW{{\mathcal W}}
\def\cX{{\mathcal X}}
\def\cY{{\mathcal Y}}
\def\cZ{{\mathcal Z}}
\def\rK{{\mathrm K}}
\def\fA{{\mathfrak A}}
\def\fB{{\mathfrak B}}
\def\fC{{\mathfrak C}}
\def\fD{{\mathfrak D}}
\def\fE{{\mathfrak E}}
\def\fF{{\mathfrak F}}
\def\fG{{\mathfrak G}}
\def\fH{{\mathfrak H}}
\def\fI{{\mathfrak I}}
\def\fJ{{\mathfrak J}}
\def\fK{{\mathfrak K}}
\def\fL{{\mathfrak L}}
\def\fM{{\mathfrak M}}
\def\fN{{\mathfrak N}}
\def\fO{{\mathfrak O}}
\def\fP{{\mathfrak P}}
\def\fQ{{\mathfrak Q}}
\def\fR{{\mathfrak R}}
\def\fS{{\mathfrak S}}
\def\fT{{\mathfrak T}}
\def\fU{{\mathfrak U}}
\def\fV{{\mathfrak V}}
\def\fW{{\mathfrak W}}
\def\fX{{\mathfrak X}}
\def\fY{{\mathfrak Y}}
\def\fZ{{\mathfrak Z}}
\def\rK{{\mathrm K}}
\def\cF{{\overline{\mathcal M}}}
\def\oP{{\overline{P}}}
\def\md{{\mathfrak d}}
\def\mg{{\mathfrak g}}
\def\mo{{\mathfrak o}}
\def\mm{{\mathfrak m}}
\def\mp{{\mathfrak p}}
\def\mh{{\mathfrak h}}
\def\fS{{\mathfrak S}}
\def\sg{{\mathsf g}}
\def\msfE{{\mathsf E}}
\def\bA{{\mathbb A}}
\def\bG{{\mathbb G}}
\def\bP{{\mathbb P}}
\def\bQ{{\mathbb Q}}
\def\bH{{\mathbb H}}
\def\bZ{{\mathbb Z}}
\def\bR{{\mathbb R}}
\def\bM{{\mathbb M}}
\def\bN{{\mathbb N}}
\def\bC{{\mathbb C}}
\def\F{{\mathbb F}}
\def\C{C_N\times C_{MN}}
\newcommand\Cmn[2]{C/#1\C\times\C/#2\C}
\def\lcm{\mathrm{lcm}}
\def\rH{{\mathrm H}}
\def\rK{{\mathrm K}}
\def\rI{{\mathrm I}}
\def\gm{{\mathbb G}_m}
\DeclareMathOperator{\Ima}{Im}
\DeclareMathOperator{\leng}{length}
\def\bF{{\mathbb F}}
\def\vol{\mathrm{vol}}
\def\Br{\mathrm{Br}}
\def\codim{\mathrm{codim}}
\def\res{\mathrm{res}}
\def\CH{\mathrm{CH}}
\def\BC{\mathcal{BC}}
\def\Bl{\mathrm{Bl}}
\def\Pic{\mathrm{Pic}}
\def\End{\mathrm{End}}
\def\ev{\mathrm{ev}}
\def\Gal{\mathrm{Gal}}
\def\Aut{\mathrm{Aut}}
\def\Gr{\mathrm{Gr}}
\def\GIT{/\!\!/}
\def\SL{\mathsf{SL}}
\def\GL{\mathsf{GL}}
\def\bPGL{\mathsf{PGL}}
\def\PGL{\mathsf{PGL}}
\def\OGr{\mathrm{OGr}}
\def\Ext{\mathrm{Ext}}
\def\Am{\mathrm{Am}}
\def\Hom{\mathrm{Hom}}
\def\Orb{\mathrm{Orb}}
\def\Isom{\mathrm{Isom}}
\def\Conj{\mathrm{Conj}}
\def\Burn{\mathrm{Burn}}
\def\Bir{\mathrm{Bir}}
\def\et{\mathrm{et}}
\def\lim{\mathrm{lim}}
\def\NS{\mathrm{NS}}
\def\trdeg{\mathrm{trdeg}}
\def\Prim{\mathrm{prim}}
\def\Orb{\mathrm{Orb}}
\def\Pt{\mathrm{pt.}}
\def\Prym{\mathrm{Prym}}
\def\Sect{\mathrm{Sect}}
\def\Sym{\mathrm{Sym}}
\def\Var{\mathrm{Var}}
\def\IJ{\mathrm{IJ}}
\def\JJ{\mathrm{J}}
\def\ch{\mathrm{char}}
\def\Spec{\mathrm{Spec}}
\def\Ker{\mathrm{Ker}}
\def\ori{\mathrm{or}}
\def\Cr{\mathrm{Cr}}
\def\cF{\mathcal F}
\def\cd{\mathrm{cd}}
\def\wcV{\widetilde{\cV}}
\def\wcX{\widetilde{\cX}}
\begin{document}
\title[Linearizability notions in equivariant birational geometry]{Linearizability notions in equivariant birational geometry}
\author{New York, NY 10012; USA Simons Foundation; 160 Fifth Av; Andrew Kresch; Yuri Tschinkel}
\date{}
\maketitle
\noindent\emph{Source and licence.} Linearizability notions in equivariant birational geometry. Version arXiv:2606.10965v1 (CC BY 4.0). This material is distributed under the Creative Commons Attribution 4.0 International licence, http://creativecommons.org/licenses/by/4.0/, as recorded in the arXiv OAI licence metadata for this exact version and shown on the abstract page https://arxiv.org/abs/2606.10965v1. Full rights notice: licenses/arxiv-2606.10965-CC-BY-4.0.txt.\par\smallskip
\noindent\textit{Editorial scope.} Counted unit: the original \texttt{theo} environment \texttt{thm.onetorsorP} at source lines 497--518 together with its complete proof at lines 520--545; the delivered document keeps source lines 349--546 so that every referenced label and every citation resolves inside it. Native editable source is the paper's own concatenated TeX; the AMS wrapper is editorial formatting only. No publisher class, upstream script or unrelated artwork is redistributed.
\section{Basic notions of equivariant birational geometry}
\label{sect:gen}
We work over an algebraically closed field $k$ of characteristic zero.

\subsection*{Varieties over extension fields}
For an extension field $K$ of $k$,
a variety over $K$ is a scheme of finite type over $K$ that is separated, reduced and irreducible.
A proper birational map $X\dashrightarrow Y$ of varieties over $K$ is one that can be factored into proper birational morphisms and their formal inverses.
Existence of a proper birational map over $K$ defines an equivalence relation, for which we employ the notation
\[ X\sim Y. \]
For the stable version, with $X$ replaced by $X\times \bP^m$ and $Y$ by $Y\times \bP^n$, for some $m$ and $n$, we employ the notation
\[ X\sim^s Y. \]
%Two varieties $X$ and $Y$ over $K$ are birationally equivalent if $K(X)\cong K(Y)$, and stably birationally equivalent if $X\times \bP^m$ and $Y\times\bP^n$ are birationally equivalent for some $m$ and $n$. 
We recall well-studied conditions for a variety $X$ over $K$ of dimension $d$ (understood for general $X$ as applied to a proper model of $X$):
\begin{itemize}
\item {\em rationality}: 
$X\sim\bP^d$, 
\item {\em stable rationality}: $X\sim^s\bP^d$.
\end{itemize}




\subsection*{$G$-varieties}
Let $G$ be a finite group.
A $G$-variety (over $k$) is a separated, reduced, finite-type $k$-scheme $X$ with regular $G$-action, such that the $G$-action is transitive on the set of irreducible components of $X$.
We call $X$ an irreducible $G$-variety if, non-equivariantly, $X$ is irreducible, i.e., a variety.
Given a $G$-variety $X$, the $G$-action is generically free if it is free on an invariant dense open subvariety of $X$.
By convention, the actions are right $G$-actions.

We recall some further notions; see, e.g., \cite[\S 3.1]{KTsurvey} for more details.
A $G$-rational map is a $G$-equivariant rational map between $G$-varieties.
A $G$-birational map is a $G$-rational map that is birational.
A proper $G$-birational map is one that factors as a composition of proper birational $G$-equivariant morphisms and their formal inverses.





Among $G$-varieties we have the relation, to be \emph{$G$-birational}:
\[ X\sim_GY\quad\Leftrightarrow\quad
\exists\text{ proper $G$-birational map $X\dashrightarrow Y$}. \]
We call $X$ and $Y$ \emph{stably $G$-birational} if they become $G$-birational after passing to products with (nonempty) projectivized representations:
\[ X\sim_G^sY\quad \Leftrightarrow\quad
\exists\text{ $G$-representations $U$, $V$}\colon X\times \bP(U)\sim_G Y\times \bP(V). \]
We also recall birationally invariant conditions on a $G$-variety $X$:
\begin{itemize}
\item \textbf{(L)} \emph{linearizability}: $X\sim_G\bP(V)$ for some representation $V$.
\item \textbf{(SL)} \emph{stable linearizability} $X\times \bP(U)\sim_G\bP(V)$ for some representations $U$ and $V$.
\item \textbf{(U)} \emph{$G$-unirationality} if there is a dominant $G$-rational map $\bP(V)\dashrightarrow X$ for some representation $V$.
\end{itemize}
(These conditions apply as stated when $X$ is proper and are understood by convention for more general $X$ when they are satisfied by an equivariant proper model; cf.\ \cite[\S 3.1]{KTsurvey}.)

\subsection*{Torsors}
Further notions are defined in terms of torsors, by which we mean torsors under $G$.
Without further mention, the base of the torsor will be $\Spec(K)$, where $K$ is an extension field of $k$.
The total space will be $T=\Spec(L)$ for an \'etale $K$-algebra $L$.
The next condition on a $G$-variety $X$ is defined by requiring a condition to hold for all torsors, over all extension fields $K$ of $k$ (expressed simply as $\forall\,T\colon\dots$).
\begin{itemize}
\item \textbf{(WV)} $X$ is \emph{weakly versal} if $\forall\,T$: $\exists$ $G$-equivariant $\Spec(L)\to X$.
\item \textbf{(V)} $X$ is \emph{versal} if every invariant dense open subvariety of $X$ is weakly versal.
\end{itemize}

Let $Z$ be a quasi-projective $G$-variety over $k$ with a generically free action of $G$ and $Z\to Z/G$ the quotient.
The function field $k(Z/G)$ is the field of $G$-invariants $k(Z)^G$.
We get an associated torsor
\[ T_Z:=k(Z)/k(Z)^G, \]
corresponding to the Galois field extension,
with Galois group $G$.
An important special case $T_V$ arises from a faithful $G$-representation $V$ and the quotient $V\to V/G$. 



\subsection*{Twisting}
Given a torsor $T=\Spec(L)$ over $K$, we obtain a free $G$-action on the $G$-variety $$
X_L=X\times_{\Spec(k)}\Spec(L)
$$ 
over $K$.
The quotient variety for this free $G$-action, provided it exists, is the \emph{twist} ${}^TX$,
a variety over $K$ with the property
$({}^TX)_L\cong X_L$.
The quotient variety is guaranteed to exist when $X$ is quasi-projective.
In the sequel we will apply the twisting construction only to quasi-projective varieties.

When applied to a torsor $T_Z$ arising from a generically free 
$G$-action on $Z$ as above, we get the twist ${}^{T_Z}X$ over $K=k(Z)^G$.

\begin{exam}
\label{exa.Vtwist}
For a $G$-representation $V$ of dimension $n\ge 1$, we have
${}^TV\cong \bA^n$ \cite[Lemma 3.1]{DR} and
${}^T\bP(V)\cong \bP^{n-1}$, for all torsors $T$.
\end{exam}

\begin{rema}
    \label{rema:twist-stab}
The twisting construction respects $G$-invariant loci $U\subset X$, i.e.,
${}^TU\subset {}^TX$, for all $G$-torsors $T$. In particular, it preserves the stabilizer stratification on $X$, see \cite[Sect.\ 6.1]{KTsurvey}.
\end{rema}









\subsection*{Versality and twisting}



Duncan and Reichstein give the following characterizations.

\begin{prop}[{\cite{DR}}]
\label{prop.DR}
Let $X$ be a quasi-projective $G$-variety.
\begin{itemize}
\item[(i)] $X$ is weakly versal $\Leftrightarrow\,\forall\,T\colon{}^TX(K)\ne \emptyset$.
\item[(ii)] $X$ is versal $\Leftrightarrow\,\forall\,T\colon{}^TX(K)$ is dense in ${}^TX$.
\item[(iii)] $X$ is $G$-unirational $\Leftrightarrow\,\forall\,T\colon{}^TX$ is unirational.
\item[(iv)] $X$ is stably linearizable $\Leftrightarrow\,\forall\,T\colon{}^TX$ is stably rational.
\end{itemize}
\end{prop}

For smooth projective $G$-varieties, these conditions are stably birationally invariant.
Generally, we have the implications \cite[(1.1)]{DR}:
\[ \textbf{(SL)}\quad\Rightarrow\quad \textbf{(U)}\quad\Rightarrow\quad \textbf{(V)}\quad\Rightarrow\quad \textbf{(WV)}. \]
(In \cite{DR}, group actions satisfying \textbf{(U)} are called \emph{very versal}.) 
There is also the implication
\[ \textbf{(WV)}\quad\Rightarrow\quad \textbf{(A)} \]
\cite[Rmk.\ 2.7]{DR},
where Condition \textbf{(A)} is the existence of fixed points for all abelian subgroups of $G$.

\begin{exam}
\label{exa.WVV}
All of these implications are strict.
For the first two, we see this already when
$G$ is trivial \cite[Sect.\ 1]{DR}.
For $G=\bZ/2\bZ$, the action on an elliptic curve $E$ by $-1$ has fixed points, thus is weakly versal \cite[Prop.\ 2.2]{DR}, but is not versal:
for $T\colon k(t)/k(t^2)$, the only rational points on the twist ${}^TE$ are the ones coming from the fixed points.
Actions satisfying condition \textbf{(A)}, but not \textbf{(WV)}, may be found among toric varieties and del Pezzo surfaces \cite{KT-uni}, \cite{STZ}.
\end{exam}

In the proof of their characterizations by means of torsors, Duncan and Reichstein make the observation that a property for all torsors is implied by the property for a single torsor, coming from a faithful $G$-representation.
Such observations had also appeared earlier, e.g., as recalled in \cite[Prop.\ 2.1]{GM}.
Our next result presents this observation uniformly for all of the relevant properties and in the strengthened form, that the single torsor can be given by any generically free \emph{versal} group action.


\begin{theo}
\label{thm.onetorsorP}
Let $X$ be a quasi-projective $G$-variety over $k$, and
$\mathbf{P}$ one of the following properties of a variety over an extension field of $k$: 
\begin{itemize}
    \item possession of a rational point,
\item possession of a Zariski dense set of rational points,
\item unirationality,
\item stable rationality,
\item rationality.
\end{itemize}
The following are equivalent.
\begin{itemize}
\item[(i)] For every torsor $T$, the twist
${}^TX$ satisfies property $\mathbf{P}$.
\item[(ii)] For some quasi-projective $G$-variety $Z$ over $k$ with generically free versal $G$-action and 
$T_Z=k(Z)/k(Z)^G$, 
the twist ${}^{T_Z}X$ 
%${}^{k(Z)/k(Z)^G}X$ 
satisfies property $\mathbf{P}$.
\end{itemize}
\end{theo}

\begin{proof}
The implication (i) $\Rightarrow$ (ii) is trivial.
We suppose (ii), and let
$K$ be an extension field of $k$ and
$T\colon L/K$ a $G$-torsor, where
$L$ is an \'etale $K$-algebra.
We need to establish property $\mathbf{P}$
for the twist ${}^TX$.
We first treat the case when $\mathbf{P}$ is ``possession of a rational point'', then we explain the modifications to the argument for
the other properties $\mathbf{P}$.
If ${}^{T_Z}X$ has a rational point, then there is a $G$-rational map $Z\dashrightarrow X$, say, defined on
invariant dense open $U\subset Z$.
Since the $G$-action on $Z$ is versal, there is a $G$-equivariant morphism $\Spec(L)\to U$.
Composing, we get equivariant $\Spec(L)\to X$, thus a $K$-point of ${}^TX$.

Generally, density of a given set of points of a finite-type scheme over a field is stable under field extension (reduce to the case of a finitely generated field extension and treat the finite and purely transcendental cases separately).
So, if ${}^{T_Z}X$ has a dense set of $k(Z)^G$-points, then
${}^{K(Z)/K(Z)^G}X$ has a dense set of $K(Z)^G$-points.
For invariant dense open $W\subset X_K$ we argue as above with $Z_K\dashrightarrow W$, defined on invariant dense open $U\subset Z_K$, and get a $K$-point of ${}^TX$, for which the corresponding $G$-equivariant morphism $\Spec(L)\to X_K$ has image in $W$.

If ${}^{T_Z}X$ is unirational, then for some $N$ there exists a dominant $G$-rational map $\bP^N\times Z\dashrightarrow X\times Z$ (with trivial action on $\bP^N$), compatible with the projection morphisms to $Z$.
We choose an invariant dense open subset where this is a morphism and an invariant dense open subset $W$ of the image and let $U\subset Z$ denote the image of $W\subset X\times Z$ under projection.
As before, by versality, there is equivariant $\Spec(L)\to U$, and by base change, dominant equivariant $\bP^N_L\dashrightarrow X_L$, hence dominant $\bP^N_K\dashrightarrow {}^TX$.

For the remaining properties, we get the result as a special case of the next theorem.
\end{proof}


\begin{thebibliography}{99}
\bibitem[DR15]{DR} A.~Duncan and Z.~Reichstein. \newblock Versality of algebraic group actions and rational points on twisted varieties. \newblock {\em J. Algebraic Geom.}, 24(3):499--530, 2015. \newblock With an appendix containing a letter from J.-P. Serre.
\bibitem[GM22]{GM} M.~Gherman and A.~Merkurjev. \newblock Negligible degree two cohomology of finite groups. \newblock {\em J. Algebra}, 611:82--93, 2022.
\bibitem[KT25]{KT-uni} A.~Kresch and Yu. Tschinkel. \newblock Equivariant unirationality of toric varieties, 2025. \newblock {\tt arXiv:2506.07152}.
\bibitem[KT26]{KTsurvey} A.~Kresch and Yu. Tschinkel. \newblock Invariants in equivariant birational geometry, 2026. \newblock {\tt arXiv:2602.23998}, to appear in {\em EMS Surv. Math. Sci.}
\bibitem[STZ26]{STZ} F.~Scavia, Yu. Tschinkel, and Zh. Zhang. \newblock Birational invariance of higher {A}mitsur groups, 2026. \newblock {\tt arXiv:2605.02763}.
\end{thebibliography}
\end{document}
